    
    

    
    \hypertarget{a-minimal-alf-run}{%
\section{A minimal ALF run}\label{a-minimal-alf-run}}

    In this bare-bones example we use the
\href{https://git.physik.uni-wuerzburg.de/ALF/pyALF/-/tree/ALF-2.4/}{pyALF}
interface to run the canonical Hubbard model on a default configuration:
a \(6\times6\) square grid, with interaction strength \(U=4\) and
inverse temperature \(\beta = 5\).

Bellow we go through the steps for performing the simulation and
outputting observables.

\begin{center}\rule{0.5\linewidth}{0.5pt}\end{center}

\textbf{1.} Import \texttt{ALF\_source} and \texttt{Simulation} classes
from the \texttt{py\_alf} python module, which provide the interface
with ALF:

    \begin{tcolorbox}[breakable, size=fbox, boxrule=1pt, pad at break*=1mm,colback=cellbackground, colframe=cellborder]
\prompt{In}{incolor}{1}{\boxspacing}
\begin{Verbatim}[commandchars=\\\{\}]
\PY{k+kn}{from} \PY{n+nn}{py\PYZus{}alf} \PY{k+kn}{import} \PY{n}{ALF\PYZus{}source}\PY{p}{,} \PY{n}{Simulation}  \PY{c+c1}{\PYZsh{} Interface with ALF}
\end{Verbatim}
\end{tcolorbox}

    \textbf{2.} Create an instance of \texttt{ALF\_source}, downloading the
ALF source code from the
\href{https://git.physik.uni-wuerzburg.de/ALF/ALF/-/tree/master/}{ALF
repository}, if \texttt{alf\_dir} does not exist. Gets alf\_dir from
environment variable \texttt{\$ALF\_DIR}, or defaults to ``./ALF'', if
not present:

    \begin{tcolorbox}[breakable, size=fbox, boxrule=1pt, pad at break*=1mm,colback=cellbackground, colframe=cellborder]
\prompt{In}{incolor}{2}{\boxspacing}
\begin{Verbatim}[commandchars=\\\{\}]
\PY{n}{alf\PYZus{}src} \PY{o}{=} \PY{n}{ALF\PYZus{}source}\PY{p}{(}\PY{n}{branch}\PY{o}{=}\PY{l+s+s1}{\PYZsq{}}\PY{l+s+s1}{ALF\PYZhy{}2.4}\PY{l+s+s1}{\PYZsq{}}\PY{p}{)}
\end{Verbatim}
\end{tcolorbox}

    \begin{Verbatim}[commandchars=\\\{\}]
Checking out branch ALF-2.4
Your branch is up to date with 'origin/ALF-2.4'.
    \end{Verbatim}

    \begin{Verbatim}[commandchars=\\\{\}]
Already on 'ALF-2.4'
    \end{Verbatim}

    \textbf{3.} Create an instance of \texttt{Simulation}, overwriting
default parameters as desired:

    \begin{tcolorbox}[breakable, size=fbox, boxrule=1pt, pad at break*=1mm,colback=cellbackground, colframe=cellborder]
\prompt{In}{incolor}{3}{\boxspacing}
\begin{Verbatim}[commandchars=\\\{\}]
\PY{n}{sim} \PY{o}{=} \PY{n}{Simulation}\PY{p}{(}
    \PY{n}{alf\PYZus{}src}\PY{p}{,}
    \PY{l+s+s2}{\PYZdq{}}\PY{l+s+s2}{Hubbard}\PY{l+s+s2}{\PYZdq{}}\PY{p}{,}                    \PY{c+c1}{\PYZsh{} Name of Hamiltonian}
    \PY{p}{\PYZob{}}                             \PY{c+c1}{\PYZsh{} Dictionary overwriting default parameters}
        \PY{l+s+s2}{\PYZdq{}}\PY{l+s+s2}{Lattice\PYZus{}type}\PY{l+s+s2}{\PYZdq{}}\PY{p}{:} \PY{l+s+s2}{\PYZdq{}}\PY{l+s+s2}{Square}\PY{l+s+s2}{\PYZdq{}}
    \PY{p}{\PYZcb{}}\PY{p}{,}
    \PY{n}{machine}\PY{o}{=}\PY{l+s+s1}{\PYZsq{}}\PY{l+s+s1}{GNU}\PY{l+s+s1}{\PYZsq{}}  \PY{c+c1}{\PYZsh{} Change to \PYZdq{}intel\PYZdq{}, or \PYZdq{}PGI\PYZdq{} if gfortran is not installed}
\PY{p}{)}
\end{Verbatim}
\end{tcolorbox}

    \textbf{4.} Compile ALF. The first time it will also download and
compile HDF5, which could take \(\sim15\) minutes.

    \begin{tcolorbox}[breakable, size=fbox, boxrule=1pt, pad at break*=1mm,colback=cellbackground, colframe=cellborder]
\prompt{In}{incolor}{4}{\boxspacing}
\begin{Verbatim}[commandchars=\\\{\}]
\PY{n}{sim}\PY{o}{.}\PY{n}{compile}\PY{p}{(}\PY{p}{)}
\end{Verbatim}
\end{tcolorbox}

    \begin{Verbatim}[commandchars=\\\{\}]
Compiling ALF{\ldots}
Cleaning up Prog/
Cleaning up Libraries/
Cleaning up Analysis/
Compiling Libraries
    \end{Verbatim}

    \begin{Verbatim}[commandchars=\\\{\}]
entanglement\_mod.F90:35:2:

   35 | \#warning "You are compiling entanglement without MPI. No results
possible"
      |  1\textasciitilde{}\textasciitilde{}\textasciitilde{}\textasciitilde{}\textasciitilde{}\textasciitilde{}
Warning: \#warning "You are compiling entanglement without MPI. No results
possible" [-Wcpp]
ar: creating modules\_90.a
ar: creating libqrref.a
    \end{Verbatim}

    \begin{Verbatim}[commandchars=\\\{\}]
Compiling Analysis
Compiling Program
Parsing Hamiltonian parameters
filename: Hamiltonians/Hamiltonian\_Kondo\_smod.F90
filename: Hamiltonians/Hamiltonian\_Hubbard\_smod.F90
filename: Hamiltonians/Hamiltonian\_Hubbard\_Plain\_Vanilla\_smod.F90
filename: Hamiltonians/Hamiltonian\_tV\_smod.F90
filename: Hamiltonians/Hamiltonian\_LRC\_smod.F90
filename: Hamiltonians/Hamiltonian\_Z2\_Matter\_smod.F90
Compiling program modules
Link program
Done.
    \end{Verbatim}

    \textbf{5.} Perform the simulation as specified in \texttt{sim}:

    \begin{tcolorbox}[breakable, size=fbox, boxrule=1pt, pad at break*=1mm,colback=cellbackground, colframe=cellborder]
\prompt{In}{incolor}{5}{\boxspacing}
\begin{Verbatim}[commandchars=\\\{\}]
\PY{n}{sim}\PY{o}{.}\PY{n}{run}\PY{p}{(}\PY{p}{)}
\end{Verbatim}
\end{tcolorbox}

    \begin{Verbatim}[commandchars=\\\{\}]
Prepare directory "/home/jonas/Programs/pyALF/Notebooks/ALF\_data/Hubbard\_Square"
for Monte Carlo run.
Create new directory.
Run /home/jonas/Programs/ALF/Prog/ALF.out
 ALF Copyright (C) 2016 - 2021 The ALF project contributors
 This Program comes with ABSOLUTELY NO WARRANTY; for details see license.GPL
 This is free software, and you are welcome to redistribute it under certain
conditions.
 No initial configuration
    \end{Verbatim}

    \textbf{6.} Perform some simple analysis:

    \begin{tcolorbox}[breakable, size=fbox, boxrule=1pt, pad at break*=1mm,colback=cellbackground, colframe=cellborder]
\prompt{In}{incolor}{6}{\boxspacing}
\begin{Verbatim}[commandchars=\\\{\}]
\PY{n}{sim}\PY{o}{.}\PY{n}{analysis}\PY{p}{(}\PY{p}{)}
\end{Verbatim}
\end{tcolorbox}

    \begin{Verbatim}[commandchars=\\\{\}]
\#\#\# Analyzing /home/jonas/Programs/pyALF/Notebooks/ALF\_data/Hubbard\_Square \#\#\#
/home/jonas/Programs/pyALF/Notebooks
Scalar observables:
Ener\_scal
Kin\_scal
Part\_scal
Pot\_scal
Histogram observables:
Equal time observables:
Den\_eq
Green\_eq
SpinT\_eq
SpinXY\_eq
SpinZ\_eq
Time displaced observables:
Den\_tau
Green\_tau
SpinT\_tau
SpinXY\_tau
SpinZ\_tau
    \end{Verbatim}

    \textbf{7.} Read analysis results into a Pandas Dataframe with one row
per simulation, containing parameters and observables:

    \begin{tcolorbox}[breakable, size=fbox, boxrule=1pt, pad at break*=1mm,colback=cellbackground, colframe=cellborder]
\prompt{In}{incolor}{7}{\boxspacing}
\begin{Verbatim}[commandchars=\\\{\}]
\PY{n}{obs} \PY{o}{=} \PY{n}{sim}\PY{o}{.}\PY{n}{get\PYZus{}obs}\PY{p}{(}\PY{p}{)}
\end{Verbatim}
\end{tcolorbox}

    \begin{Verbatim}[commandchars=\\\{\}]
/home/jonas/Programs/pyALF/Notebooks/ALF\_data/Hubbard\_Square
No orbital locations saved.
    \end{Verbatim}

    \begin{tcolorbox}[breakable, size=fbox, boxrule=1pt, pad at break*=1mm,colback=cellbackground, colframe=cellborder]
\prompt{In}{incolor}{8}{\boxspacing}
\begin{Verbatim}[commandchars=\\\{\}]
\PY{n}{obs}
\end{Verbatim}
\end{tcolorbox}

            \begin{tcolorbox}[breakable, size=fbox, boxrule=.5pt, pad at break*=1mm, opacityfill=0]
\prompt{Out}{outcolor}{8}{\boxspacing}
\begin{Verbatim}[commandchars=\\\{\}]
                                                    continuous  ham\_chem  \textbackslash{}
/home/jonas/Programs/pyALF/Notebooks/ALF\_data/H{\ldots}           0       0.0

                                                    ham\_t  ham\_t2  ham\_tperp  \textbackslash{}
/home/jonas/Programs/pyALF/Notebooks/ALF\_data/H{\ldots}    1.0     1.0        1.0

                                                    ham\_u  ham\_u2  mz  l1  l2  \textbackslash{}
/home/jonas/Programs/pyALF/Notebooks/ALF\_data/H{\ldots}    4.0     4.0   1   6   6

                                                    {\ldots}  \textbackslash{}
/home/jonas/Programs/pyALF/Notebooks/ALF\_data/H{\ldots}  {\ldots}

      SpinXY\_tauK\_err  \textbackslash{}
/home/jonas/Programs/pyALF/Notebooks/ALF\_data/H{\ldots}  [[0.22628722504550844,
0.37685459854025805, 0{\ldots}

          SpinXY\_tauR  \textbackslash{}
/home/jonas/Programs/pyALF/Notebooks/ALF\_data/H{\ldots}  [[0.0575988003757706,
-0.10378742200171556, 0{\ldots}

      SpinXY\_tauR\_err  \textbackslash{}
/home/jonas/Programs/pyALF/Notebooks/ALF\_data/H{\ldots}  [[0.01248834666799511,
0.036871289378260834, 0{\ldots}

   SpinXY\_tau\_lattice  \textbackslash{}
/home/jonas/Programs/pyALF/Notebooks/ALF\_data/H{\ldots}  \{'L1': [6.0, 0.0], 'L2':
[0.0, 6.0], 'a1': [1{\ldots}

           SpinZ\_tauK  \textbackslash{}
/home/jonas/Programs/pyALF/Notebooks/ALF\_data/H{\ldots}  [[0.9040036362034092,
0.5628191896020613, 0.61{\ldots}

       SpinZ\_tauK\_err  \textbackslash{}
/home/jonas/Programs/pyALF/Notebooks/ALF\_data/H{\ldots}  [[0.1494786750311893,
0.06006336638595281, 0.0{\ldots}

           SpinZ\_tauR  \textbackslash{}
/home/jonas/Programs/pyALF/Notebooks/ALF\_data/H{\ldots}  [[0.1013442053139478,
-0.11445552391592478, 0{\ldots}

       SpinZ\_tauR\_err  \textbackslash{}
/home/jonas/Programs/pyALF/Notebooks/ALF\_data/H{\ldots}  [[0.06273438448505823,
0.056390780430637014, 0{\ldots}

    SpinZ\_tau\_lattice  \textbackslash{}
/home/jonas/Programs/pyALF/Notebooks/ALF\_data/H{\ldots}  \{'L1': [6.0, 0.0], 'L2':
[0.0, 6.0], 'a1': [1{\ldots}

              lattice
/home/jonas/Programs/pyALF/Notebooks/ALF\_data/H{\ldots}  \{'L1': [6.0, 0.0], 'L2':
[0.0, 6.0], 'N\_coord'{\ldots}

[1 rows x 111 columns]
\end{Verbatim}
\end{tcolorbox}
        
    \(\bullet\) The internal energy of the system (and its error) are
accessed by:

    \begin{tcolorbox}[breakable, size=fbox, boxrule=1pt, pad at break*=1mm,colback=cellbackground, colframe=cellborder]
\prompt{In}{incolor}{9}{\boxspacing}
\begin{Verbatim}[commandchars=\\\{\}]
\PY{n}{obs}\PY{o}{.}\PY{n}{iloc}\PY{p}{[}\PY{l+m+mi}{0}\PY{p}{]}\PY{p}{[}\PY{p}{[}\PY{l+s+s1}{\PYZsq{}}\PY{l+s+s1}{Ener\PYZus{}scal0}\PY{l+s+s1}{\PYZsq{}}\PY{p}{,} \PY{l+s+s1}{\PYZsq{}}\PY{l+s+s1}{Ener\PYZus{}scal0\PYZus{}err}\PY{l+s+s1}{\PYZsq{}}\PY{p}{,} \PY{l+s+s1}{\PYZsq{}}\PY{l+s+s1}{Ener\PYZus{}scal\PYZus{}sign}\PY{l+s+s1}{\PYZsq{}}\PY{p}{,} \PY{l+s+s1}{\PYZsq{}}\PY{l+s+s1}{Ener\PYZus{}scal\PYZus{}sign\PYZus{}err}\PY{l+s+s1}{\PYZsq{}}\PY{p}{]}\PY{p}{]}
\end{Verbatim}
\end{tcolorbox}

            \begin{tcolorbox}[breakable, size=fbox, boxrule=.5pt, pad at break*=1mm, opacityfill=0]
\prompt{Out}{outcolor}{9}{\boxspacing}
\begin{Verbatim}[commandchars=\\\{\}]
Ener\_scal0           -29.821914
Ener\_scal0\_err          0.13032
Ener\_scal\_sign              1.0
Ener\_scal\_sign\_err          0.0
Name: /home/jonas/Programs/pyALF/Notebooks/ALF\_data/Hubbard\_Square, dtype:
object
\end{Verbatim}
\end{tcolorbox}
        
    \(\bullet\) The simulation can be resumed by calling \texttt{sim.run()}
again, increasing the precision of results:

    \begin{tcolorbox}[breakable, size=fbox, boxrule=1pt, pad at break*=1mm,colback=cellbackground, colframe=cellborder]
\prompt{In}{incolor}{10}{\boxspacing}
\begin{Verbatim}[commandchars=\\\{\}]
\PY{n}{sim}\PY{o}{.}\PY{n}{run}\PY{p}{(}\PY{p}{)}
\PY{n}{sim}\PY{o}{.}\PY{n}{analysis}\PY{p}{(}\PY{p}{)}
\PY{n}{obs2} \PY{o}{=} \PY{n}{sim}\PY{o}{.}\PY{n}{get\PYZus{}obs}\PY{p}{(}\PY{p}{)}
\PY{n}{obs2}\PY{o}{.}\PY{n}{iloc}\PY{p}{[}\PY{l+m+mi}{0}\PY{p}{]}\PY{p}{[}\PY{p}{[}\PY{l+s+s1}{\PYZsq{}}\PY{l+s+s1}{Ener\PYZus{}scal0}\PY{l+s+s1}{\PYZsq{}}\PY{p}{,} \PY{l+s+s1}{\PYZsq{}}\PY{l+s+s1}{Ener\PYZus{}scal0\PYZus{}err}\PY{l+s+s1}{\PYZsq{}}\PY{p}{,} \PY{l+s+s1}{\PYZsq{}}\PY{l+s+s1}{Ener\PYZus{}scal\PYZus{}sign}\PY{l+s+s1}{\PYZsq{}}\PY{p}{,} \PY{l+s+s1}{\PYZsq{}}\PY{l+s+s1}{Ener\PYZus{}scal\PYZus{}sign\PYZus{}err}\PY{l+s+s1}{\PYZsq{}}\PY{p}{]}\PY{p}{]}
\end{Verbatim}
\end{tcolorbox}

    \begin{Verbatim}[commandchars=\\\{\}]
Prepare directory "/home/jonas/Programs/pyALF/Notebooks/ALF\_data/Hubbard\_Square"
for Monte Carlo run.
Resuming previous run.
Run /home/jonas/Programs/ALF/Prog/ALF.out
 ALF Copyright (C) 2016 - 2021 The ALF project contributors
 This Program comes with ABSOLUTELY NO WARRANTY; for details see license.GPL
 This is free software, and you are welcome to redistribute it under certain
conditions.
\#\#\# Analyzing /home/jonas/Programs/pyALF/Notebooks/ALF\_data/Hubbard\_Square \#\#\#
/home/jonas/Programs/pyALF/Notebooks
Scalar observables:
Ener\_scal
Kin\_scal
Part\_scal
Pot\_scal
Histogram observables:
Equal time observables:
Den\_eq
Green\_eq
SpinT\_eq
SpinXY\_eq
SpinZ\_eq
Time displaced observables:
Den\_tau
Green\_tau
SpinT\_tau
SpinXY\_tau
SpinZ\_tau
/home/jonas/Programs/pyALF/Notebooks/ALF\_data/Hubbard\_Square
No orbital locations saved.
    \end{Verbatim}

            \begin{tcolorbox}[breakable, size=fbox, boxrule=.5pt, pad at break*=1mm, opacityfill=0]
\prompt{Out}{outcolor}{10}{\boxspacing}
\begin{Verbatim}[commandchars=\\\{\}]
Ener\_scal0           -29.609245
Ener\_scal0\_err         0.136803
Ener\_scal\_sign              1.0
Ener\_scal\_sign\_err          0.0
Name: /home/jonas/Programs/pyALF/Notebooks/ALF\_data/Hubbard\_Square, dtype:
object
\end{Verbatim}
\end{tcolorbox}
        
    \begin{tcolorbox}[breakable, size=fbox, boxrule=1pt, pad at break*=1mm,colback=cellbackground, colframe=cellborder]
\prompt{In}{incolor}{11}{\boxspacing}
\begin{Verbatim}[commandchars=\\\{\}]
\PY{n+nb}{print}\PY{p}{(}\PY{l+s+s2}{\PYZdq{}}\PY{l+s+se}{\PYZbs{}n}\PY{l+s+s2}{Running again reduced the error from}\PY{l+s+se}{\PYZbs{}n}\PY{l+s+s2}{\PYZdq{}}\PY{p}{,}
      \PY{n}{obs}\PY{o}{.}\PY{n}{iloc}\PY{p}{[}\PY{l+m+mi}{0}\PY{p}{]}\PY{p}{[}\PY{p}{[}\PY{l+s+s1}{\PYZsq{}}\PY{l+s+s1}{Ener\PYZus{}scal0\PYZus{}err}\PY{l+s+s1}{\PYZsq{}}\PY{p}{]}\PY{p}{]}\PY{p}{,} \PY{l+s+s2}{\PYZdq{}}\PY{l+s+se}{\PYZbs{}n}\PY{l+s+s2}{to}\PY{l+s+se}{\PYZbs{}n}\PY{l+s+s2}{\PYZdq{}}\PY{p}{,}
      \PY{n}{obs2}\PY{o}{.}\PY{n}{iloc}\PY{p}{[}\PY{l+m+mi}{0}\PY{p}{]}\PY{p}{[}\PY{p}{[}\PY{l+s+s1}{\PYZsq{}}\PY{l+s+s1}{Ener\PYZus{}scal0\PYZus{}err}\PY{l+s+s1}{\PYZsq{}}\PY{p}{]}\PY{p}{]}\PY{p}{)}
\end{Verbatim}
\end{tcolorbox}

    \begin{Verbatim}[commandchars=\\\{\}]

Running again reduced the error from
 Ener\_scal0\_err    0.13032
Name: /home/jonas/Programs/pyALF/Notebooks/ALF\_data/Hubbard\_Square, dtype:
object
to
 Ener\_scal0\_err    0.136803
Name: /home/jonas/Programs/pyALF/Notebooks/ALF\_data/Hubbard\_Square, dtype:
object
    \end{Verbatim}

    \textbf{Note}: To run a fresh simulation - instead of performing a
refinement over previous run(s) - the Monte Carlo run directory should
be deleted before rerunning.

    \begin{center}\rule{0.5\linewidth}{0.5pt}\end{center}

\hypertarget{exercises}{%
\subsection{Exercises}\label{exercises}}

\begin{enumerate}
\def\labelenumi{\arabic{enumi}.}
\tightlist
\item
  Check info/help commands such as \texttt{help(alf\_src)},
  \texttt{alf\_src.get\_ham\_names()}, and
  \texttt{alf\_src.get\_default\_params(\textquotesingle{}Hubbard\textquotesingle{})}.
\item
  Rerun once again and check the new improvement in precision.
\item
  Look at a few other observables (\texttt{sim.analysis()} outputs the
  names of those available).
\item
  Change the lattice size by adding, e.g.,
  \texttt{\textbackslash{}"L1\textbackslash{}":\ 4,} and
  \texttt{\textbackslash{}"L2\textbackslash{}":\ 1,} to the simulation
  parameters definitions of \texttt{sim} (step 2).
\end{enumerate}


    % Add a bibliography block to the postdoc
    
    
    
